\section{问题分析}

\subsection{数据字段与问题类型}
题目数据同时包含类别变量和连续变量。玻璃类型、纹饰、颜色及表面风化属于类别变量；14种氧化物含量属于连续比例变量，且各成分之和应接近\(100\%\)，具有典型成分数据特征。

\begin{table}[H]
  \centering
  \caption{各问题的类型与候选方法}
  \label{tab:problem-type}
  \begin{tabularx}{\textwidth}{C{1.5cm} L{3.5cm} L{5.2cm} X}
    \toprule
    问题 & 核心任务 & 定性筛选后的候选方法 & 选择依据 \\
    \midrule
    问题一 & 类别关联、分组差异、连续值恢复 & 列联分析、Fisher/卡方检验、二元逻辑回归；非参数检验；比例修正或回归 & 风化为二分类变量，化学成分为连续变量 \\
    问题二 & 二分类与亚类划分 & 信息增益决策树、PLS--DA与分类型K-means & 大类标签已知，需给出可解释规则并识别组内组成型 \\
    问题三 & 未知样品鉴别 & 复用问题二分类器，并进行概率或边界扰动分析 & 输出为高钾/铅钡二分类 \\
    问题四 & 成分关联及类别差异 & CLR/ILR变换、Spearman/Pearson相关、差异相关分析 & 无单一因变量，且存在闭合效应 \\
    \bottomrule
  \end{tabularx}
\end{table}

\subsection{总体建模思路}
总体流程如图\ref{fig:workflow}所示。首先依据成分和阈值筛除无效样本并处理未检出值；其次研究风化与文物属性及化学成分的关系；随后建立分类和聚类模型；最后对未知样品完成判别，并比较两种玻璃的成分关联结构。

\begin{figure}[H]
  \centering
  \begin{tikzpicture}[
    >=Latex,
    every node/.style={font=\small, align=center},
    data/.style={draw=black!55, rounded corners=2pt, fill=black!6, minimum width=4.0cm, minimum height=0.85cm},
    process/.style={draw=blue!55!black, rounded corners=2pt, fill=blue!7, minimum width=4.0cm, minimum height=0.92cm},
    task/.style={draw=orange!70!black, rounded corners=2pt, fill=orange!10, minimum width=3.75cm, minimum height=1.00cm},
    arrow/.style={->, line width=0.75pt, draw=black!65}
  ]
    \node[data] (source) {附件数据\\文物属性、采样点与未知样品};
    \node[process, below=6mm of source] (preprocess) {有效性筛选与成分数据处理\\闭合、近零替换、CLR变换};
    \node[task, below left=9mm and 7mm of preprocess] (p1) {问题一\\风化规律与成分恢复};
    \node[task, below right=9mm and 7mm of preprocess] (p2) {问题二\\类型判别与亚类划分};
    \node[task, below=9mm of p1] (p4) {问题四\\两类玻璃的关联结构比较};
    \node[task, below=9mm of p2] (p3) {问题三\\未知样品鉴别与不确定性分析};
    \draw[arrow] (source) -- (preprocess);
    \draw[arrow] (preprocess) -- (p1);
    \draw[arrow] (preprocess) -- (p2);
    \draw[arrow] (p1) -- (p4);
    \draw[arrow] (p2) -- (p3);
    \draw[arrow] (p2.west) to[out=180,in=0] node[above, font=\scriptsize]{分类与亚类参照} (p4.east);
  \end{tikzpicture}
  \caption{总体建模流程}
  \label{fig:workflow}
\end{figure}

\subsection{技术难点}
\begin{enumerate}
  \item 表单2按采样点记录，同一文物可能含多个部位或不同风化程度的采样点，样本并非完全独立。
  \item 空白表示未检测到，而非普通意义上的随机缺失；对数比变换前不能直接保留零值。
  \item 各成分总和受\(100\%\)闭合约束，直接计算普通相关系数可能产生伪相关。
  \item 样本规模较小且类别不平衡，复杂模型容易过拟合，应优先选择可解释、低复杂度方法，并使用交叉验证与敏感性分析。
\end{enumerate}
